% ── SECTION 4: THE RELATIONAL AFFINITY GATE (~1,000 words) ──────────────────

\section{The Relational Affinity Gate}
\label{sec:gate}

\subsection{Definition}

The \RCT~establishes that relational affinity~$\bar\alpha$ is the fundamental
dynamical quantity governing coherence and entanglement.
A gate that implements controlled entanglement should therefore be parameterized
directly by~$\bar\alpha$, rather than operating at a fixed entanglement
capability as conventional two-qubit gates do.

\begin{definition}[Relational Affinity Gate]
The \emph{Relational Affinity Gate} is the two-qubit unitary
\begin{equation}
  \URA(\alpha) = \exp\!\left(-\frac{i\alpha}{2}
    \bigl(\sigma_x \otimes \sigma_x
         + \sigma_y \otimes \sigma_y
         + \sigma_z \otimes \sigma_z\bigr)\right),
  \label{eq:URA}
\end{equation}
where $\alpha \in [0, \pi/2]$ is the Affinity Tensor parameter and
$\sigma_{x,y,z}$ are the Pauli matrices acting on $\mathbb{C}^2$.
\end{definition}

The generator $\hat{H} = \tfrac{1}{2}(\sigma_x\otimes\sigma_x +
\sigma_y\otimes\sigma_y + \sigma_z\otimes\sigma_z)$ is the isotropic
Heisenberg exchange interaction, which generates entanglement symmetrically
across all three Pauli axes.
The parameter~$\alpha$ plays the role of an interaction time or coupling
strength, and is identified with the scalar Affinity Tensor value
$\bar\alpha$ of the two-qubit system.

In the computational basis $\{\ket{00}, \ket{01}, \ket{10}, \ket{11}\}$,
the gate takes the explicit form
\begin{equation}
  \URA(\alpha) =
  \begin{pmatrix}
    e^{-i\alpha/2}\cos(\alpha/2) + i\sin(\alpha/2) & 0 & 0 & -i e^{-i\alpha/2}\sin(\alpha/2) \\
    0 & \cos\alpha & -i\sin\alpha & 0 \\
    0 & -i\sin\alpha & \cos\alpha & 0 \\
    -i e^{-i\alpha/2}\sin(\alpha/2) & 0 & 0 & e^{-i\alpha/2}\cos(\alpha/2) + i\sin(\alpha/2)
  \end{pmatrix}.
  \label{eq:URA-matrix}
\end{equation}
At $\alpha = 0$, $\URA = \mathbf{I}_4$ (identity).
At $\alpha = \pi/4$, the gate achieves maximum entanglement capability
(see Section~\ref{subsec:weyl}).
At $\alpha = \pi/2$, $\URA$ reduces to the \textsc{swap} gate up to
local operations.

\subsection{Bell State Production}
\label{subsec:bell}

The most direct illustration of the gate's relational character is its
action on the computational basis state $\ket{00}$ at $\alpha = \pi/4$.
Applying Eq.~\eqref{eq:URA} with $\alpha = \pi/4$ and using the Hadamard
identity $\ket{00} = \tfrac{1}{\sqrt{2}}(\ket{\Phi^+} + \ket{\Phi^-})$
as an intermediate step:
\begin{equation}
  \URA(\pi/4)\ket{00}
    = \frac{1}{\sqrt{2}}\bigl(\ket{00} + \ket{11}\bigr)
    = \ket{\Phi^+},
  \label{eq:bell-production}
\end{equation}
the maximally entangled Bell state with concurrence $\mathcal{C} = 1$
and scalar affinity $\bar\alpha = 1$.
This confirms the Coherence--Affinity Correspondence (Proposition~\ref{prop:correspondence}):
the gate parameter~$\alpha$ and the resulting relational affinity~$\bar\alpha$
coincide at $\alpha = \pi/4$.

All four Bell states are generated from the two-qubit computational basis
by choosing appropriate input states and values of~$\alpha$, confirming
that $\URA$ spans the full entanglement capability of the two-qubit Hilbert
space.

\subsection{Weyl Chamber Analysis}
\label{subsec:weyl}

Two-qubit gates are classified up to local operations by their position in
the Weyl chamber of $\mathrm{SU}(4)/[\mathrm{SU}(2)\otimes\mathrm{SU}(2)]$,
parameterized by coordinates $(c_1, c_2, c_3) \in [0,\pi/2]^3$ with
$c_1 \geq c_2 \geq c_3 \geq 0$~\cite{Zhang2003WeylChamber}.
These coordinates are extracted from the \emph{local invariants} of the
two-qubit gate matrix.

For $\URA(\alpha)$, the Weyl chamber coordinates are:
\begin{equation}
  c_1(\alpha) = \alpha, \qquad
  c_2(\alpha) = \alpha, \qquad
  c_3(\alpha) = \alpha.
  \label{eq:weyl-coords}
\end{equation}
The trajectory $(c_1, c_2, c_3) = (\alpha, \alpha, \alpha)$ lies along the
main diagonal of the Weyl chamber, tracing the \emph{iSWAP family} from the
identity ($\alpha = 0$, corner $O$) to the \textsc{swap} point ($\alpha = \pi/2$).
The maximally entangling point $\alpha = \pi/4$ corresponds to the
\textsc{iswap} gate position.

This trajectory is significant: fixed standard gates occupy single
\emph{points} in the Weyl chamber~---~\textsc{cnot} at $({\pi}/{2}, 0, 0)$,
\textsc{iswap} at $(\pi/4, \pi/4, 0)$, \textsc{cz} at $(\pi/2, 0, 0)$ (locally
equivalent to \textsc{cnot}).
The gate $\URA(\alpha)$ instead traces a \emph{continuous path} parameterized
by~$\alpha$, covering a range of entanglement capabilities rather than a
single fixed point.

\subsection{Entanglement Capability}
\label{subsec:entcap}

The \emph{entanglement capability} $\mathcal{E}(\alpha)$ of a two-qubit gate
is the maximum concurrence achievable over all product input
states~\cite{Zhang2003WeylChamber}:
\begin{equation}
  \mathcal{E}(\alpha) = \max_{\ket{\phi_1}\otimes\ket{\phi_2}}
    \mathcal{C}\!\left(\URA(\alpha)\ket{\phi_1\phi_2}\right).
  \label{eq:entcap}
\end{equation}
For $\URA(\alpha)$ with Weyl coordinates $(\alpha, \alpha, \alpha)$,
the entanglement capability is:
\begin{equation}
  \mathcal{E}(\alpha) = |\sin(2\alpha)|,
  \label{eq:entcap-result}
\end{equation}
which increases monotonically from $\mathcal{E}(0) = 0$ to
$\mathcal{E}(\pi/4) = 1$ and returns to $\mathcal{E}(\pi/2) = 0$
(the \textsc{swap} gate, which permutes states without generating entanglement).

The practical implication is direct: by setting $\alpha$ to the current
$\bar\alpha$ of the device's noise environment, the gate can be operated at
the entanglement capability precisely matched to the coherence available in
hardware~---~avoiding the overreach of fixed maximally entangling gates in
high-noise regimes.
Section~\ref{sec:sim} quantifies this advantage numerically.
